\documentclass{article}
\usepackage{main}

\begin{document}
\begin{proposition}
Le nombre $\dfrac{1}{3}$ n'est pas décimal.
\end{proposition}
\begin{proof}
\begin{enumerate}
\item On annonce que l'on procède à un raisonnement par l'absurde. On suppose donc que $\dfrac{1}{3}$ \textbf{est} décimal.

\item \begin{enumerate}
\item Le cours nous dit qu'un nombre décimal est un nombre de la forme $\dfrac{a}{10^k}$. Donc il existe $a \in \Z$ et $k \in \N$ tel que
\begin{equation*}
\dfrac{1}{3} = \dfrac{a}{10^k}
\end{equation*}
\item On transforme l'égalité suivante, pour en déduire que
\begin{equation*}
10^k = 3a
\end{equation*}
\item Ainsi, \textbf{$10^k$ est un multiple de $3$}.
\end{enumerate}
\item \begin{enumerate}
\item Or, on sait que la somme des chiffres d'un multiple de $3$ est aussi un multiple de $3$.
\item Or, les chiffres de $10^k$ sont constitués d'un $1$ suivi de $k$ zéros. La somme des chiffres de $10^k$ fait donc $1$. \textbf{$10^k$ n'est pas un multiple de $3$}.
\end{enumerate} 
\item On aboutit à une \textbf{contradiction}. Par l'absurde, la supposition de départ est \textbf{fausse} : \textbf{$\dfrac{1}{3}$ n'est pas décimal}.
\end{enumerate}
\end{proof}
\end{document}